\section{What Prevents a Full-FP4 Speedup?}
\label{sec:forward-problems}

HAO's schedule already overlaps QK, softmax, and PV whenever their
dependencies allow.  Replacing its FP8 probability--value product with FP4
leaves two linked problems.  The first is timing: constructing P delays the
first useful PV instruction.  The second is numerical: softmax probabilities
span a range that is awkward for a 4-bit payload and its block scale.

\subsection{Problem 1: P lies on the critical path}

In tiled attention, P is produced inside the kernel rather than loaded as an
input.  For full FP4 attention, the first useful PV command therefore waits
for
\begin{equation}
  S\rightarrow\text{maximum}\rightarrow\text{scale}
  \rightarrow\text{probability}\rightarrow\text{E2M1 pack}
  \rightarrow\text{publish}.
  \label{eq:serial-p-path}
\end{equation}
Work removed after publication may not change latency. Work removed before
the first legal P tile can.

FP4 accelerates QK and PV, but not score reduction, exponentiation,
synchronization, or scale publication.  It also requires P to be converted
between the two products.  Removing work before the first legal P tile can
shorten the critical path; removing work after that publication may not
change latency \cite{nvidia2026ptx,zadouri2026flashattention4}.

\subsection{Problem 2: P needs both range and cheap scales}

NVFP4 and MXFP4 both use signed E2M1 payloads.  Ignoring sign, their
nonnegative magnitudes are
\begin{equation}
  \mathcal F_{\mathrm{E2M1}}=
  \left\{0,\tfrac12,1,\tfrac32,2,3,4,6\right\}.
\end{equation}
The formats differ in how they scale these payloads:
\begin{table}[H]
\centering
\caption{Block-scaled FP4 formats used in this work.}
\label{tab:fp4-formats}
\small
\begin{tabularx}{\textwidth}{lcccX}
\toprule
Format & Block & Encoded scale & Main strength & Role in this work \\
\midrule
NVFP4 & 16 & E4M3, optional tensor scale & fine local placement & signed Q and K \\
MXFP4 & 32 & E8M0 amplitude, power of two & wide exponent range & P and V operands \\
\bottomrule
\end{tabularx}
\end{table}

For an NVFP4 P block with maximum $a_B$ and global encode factor $G$, the
decoded block scale is approximately
\begin{equation}
  \widehat s_{\mathrm{NV}}=
  \frac{Q_{\mathrm{E4M3}}(G a_B/6)}{G}.
  \label{eq:nv-scale}
\end{equation}
If the E4M3 scale rounds to zero, the entire block disappears. A row-level
stabilizer fixes the range, but its scale and inverse correction add state to
the online path. MXFP4 instead stores a power-of-two block amplitude near
$a_B$,
\begin{equation}
  \alpha_B\in
  \left\{2^{\lfloor\log_2 a_B\rfloor},
          2^{\lceil\log_2 a_B\rceil}\right\},
  \qquad \delta_B=\frac{\alpha_B}{6},
  \label{eq:mx-scale}
\end{equation}
where $\alpha_B$ is the E8M0 value written to the scale page and $\delta_B$
is the effective reconstruction step.  Since the largest E2M1 code is 6, our
convention reconstructs one operand as $\alpha_B q/6$ and corrects an MXFP4
P/V product by $1/36$.  The amplitude folds naturally into a base-two score
transform and matches one 32-column producer fragment (N32). The choice is a
latency--range trade-off, not a claim that E8M0 is intrinsically more accurate
than E4M3.

Table~\ref{tab:p-range} isolates numerical range from kernel scheduling. It
forms exact Gaussian softmax probabilities, quantizes matched N32 blocks,
renormalizes the represented P, and multiplies by the same FP32 V. Stabilized
NVFP4 is the high-fidelity control; unscaled NVFP4 exposes underflow.

\begin{table}[htbp]
\centering
\caption{Probability-format range at D128. ``Zero scales'' is the fraction
of blocks whose scale encodes as zero; ``lost mass'' is exact probability
mass mapped to zero. This is a numerical diagnostic, not a kernel timing.}
\label{tab:p-range}
\scriptsize
\begin{adjustbox}{max width=\textwidth}
\begin{tabular}{llrrrrrr}
\toprule
$S$ & Format & Zero scales & Zero payload & Lost mass & $P$ rel.-$L_2$
& $PV$ cosine & $PV$ rel.-$L_2$ \\
\midrule
1024 & NVFP4 G{=}1 & 0.706299 & 0.831829 & 0.683773 & 1.715156 & 0.700819 & 1.685442 \\
1024 & NVFP4 G{=}448 & 0.000000 & 0.155418 & 0.029314 & 0.114408 & 0.993854 & 0.113687 \\
1024 & MXFP4 & 0.000000 & 0.188683 & 0.040178 & 0.146380 & 0.989846 & 0.144944 \\
4096 & NVFP4 G{=}1 & 0.996063 & 0.999438 & 0.994332 & 13.815642 & 0.217328 & 13.876892 \\
4096 & NVFP4 G{=}448 & 0.000000 & 0.156454 & 0.029760 & 0.114091 & 0.993747 & 0.114582 \\
4096 & MXFP4 & 0.000000 & 0.189755 & 0.040494 & 0.147209 & 0.989338 & 0.148716 \\
8192 & NVFP4 G{=}1 & 0.999634 & 0.999977 & 0.999284 & 12.467307 & 0.096774 & 12.389232 \\
8192 & NVFP4 G{=}448 & 0.000000 & 0.160613 & 0.030691 & 0.114725 & 0.993841 & 0.113660 \\
8192 & MXFP4 & 0.000000 & 0.194189 & 0.041630 & 0.147318 & 0.989660 & 0.146237 \\
\bottomrule
%
\end{tabular}
\end{adjustbox}
\end{table}

The alternatives therefore impose different costs.  Stabilized NVFP4 offers
the best FP4 fidelity but needs a row-level range correction.  FP8 avoids the
4-bit range problem but gives up the FP4 PV instruction.  MXFP4 provides a
power-of-two scale at the same N32 granularity as the producer fragment, but
places probabilities more coarsely.  Direct-P chooses MXFP4 and redesigns how
that operand is produced and normalized.
